\documentclass[11pt,a4paper]{article}
\usepackage{fontspec,amsmath,amssymb,amsthm,microtype}
\setmainfont{lmroman10-regular.otf}[BoldFont=lmroman10-bold.otf,ItalicFont=lmroman10-italic.otf,BoldItalicFont=lmroman10-bolditalic.otf]
\usepackage{xcolor}
\usepackage[most]{tcolorbox}
\usepackage{mathrsfs}
\providecommand{\lt}{<}
\providecommand{\gt}{>}
\setmonofont{DejaVuSansMono.ttf}[BoldFont=DejaVuSansMono-Bold.ttf,ItalicFont=DejaVuSansMono-Oblique.ttf,BoldItalicFont=DejaVuSansMono-BoldOblique.ttf,Scale=0.86]
\newfontfamily\leanfallback{STIXTwoMath-Regular.otf}[Scale=0.86]
\definecolor{leanrule}{gray}{0.62}
\definecolor{leanback}{gray}{0.962}
\definecolor{leantitle}{gray}{0.22}
\definecolor{leanmuted}{gray}{0.38}
\newtcolorbox{leanbox}[1]{enhanced,breakable,colback=leanback,colframe=leanrule,
  boxrule=0.4pt,arc=0.6mm,left=2.4mm,right=2.4mm,top=1.8mm,bottom=1.8mm,
  before skip=0.8em,after skip=1em,fonttitle=\small\bfseries,coltitle=leantitle,
  colbacktitle=leanback,titlerule=0.3pt,title={#1}}
\newlength{\leancharwidth}
\newenvironment{leancode}{\par\smallskip\begingroup\ttfamily\footnotesize\raggedright
  \settowidth{\leancharwidth}{M}\setlength{\parskip}{0pt}\setlength{\parindent}{0pt}
  \frenchspacing}{\par\endgroup\smallskip}
\newcommand{\leanline}[2]{\hangindent=\dimexpr#1\leancharwidth+4\leancharwidth\relax
  \hangafter=1\noindent\hspace*{#1\leancharwidth}#2\par}
\newcommand{\leansource}[1]{\par{\footnotesize\color{leanmuted}#1}\par}
\newcommand{\leanlabel}[1]{\par\smallskip\noindent{\small\bfseries #1}\par}
\usepackage[margin=25mm]{geometry}
\usepackage[unicode,hidelinks]{hyperref}
\hypersetup{pdfauthor={LangevinLean Project}}
\setlength{\emergencystretch}{3em}
\setlength{\parskip}{0.4em}
\setlength{\parindent}{0pt}
\allowdisplaybreaks[2]
\raggedbottom
\title{\texorpdfstring{Splitting convergence for the specified vector Brownian input}{Splitting convergence for the specified vector Brownian input}}
\author{LangevinLean Project}
\date{}
\begin{document}\maketitle
Draft mathematical exposition; these notes may contain mistakes.


The deterministic input in the splitting theorem is now identified with the specified vector Brownian motion in \(L^{2}\). Its filtered increments are exactly the constructed Gaussian segment convolutions, including the fluctuation-dissipation scaling. Consequently the innovation-driven numerical chain and the continuous-input chain coincide on the same source, and the \(L^{2}\) chain converges to its kinetic integral solution.



The continuous Brownian input, filtered increments, and full thermal stream are identified on one constructed probability source. Applying deterministic splitting convergence there gives a second-mean stochastic limit.


\section{\texorpdfstring{Keep scaling and source lifts exact}{Keep scaling and source lifts exact}}


The OU input construction is linear in the additive input. Multiplying \(B\) by a real coefficient \(c\) multiplies both its matching increment and its exponential convolution by \(c\). More generally, a continuous linear map \(L\) commutes with this construction: it commutes with subtraction, scalar multiplication, and the Bochner interval integral of the continuous integrand. One proves these exact identities by integral linearity and continuous-linear integral commutation. The latter applies, in particular, to measure-preserving pullbacks that extend the Brownian source by an independent initial state.



The Brownian source lift includes its physical amplitude before any numerical comparison is made. Since source pullback is linear, it preserves every response identity with that amplitude. This keeps thermal normalization consistent across the continuous equation and splitting increments.

\begin{leanbox}{Lean: \texttt{ouInputIncrement\_\allowbreak{}smul}, \texttt{ouInputIncrement\_\allowbreak{}map}}
\leansource{\texttt{ouInputIncrement\_\allowbreak{}smul} \quad KineticLangevin/\allowbreak{}CanonicalSplitting.lean, lines 11–22}
\begin{leancode}
\leanline{0}{theorem\ ouInputIncrement\_smul\ \{E\ :\ Type*\}\ [NormedAddCommGroup\ E]\ [NormedSpace\ ℝ\ E]}
\leanline{4}{(γ\ c\ :\ ℝ)\ (B\ :\ ℝ\ →\ E)\ (a\ h\ :\ ℝ)\ :}
\leanline{4}{ouInputIncrement\ γ\ (fun\ t\ ↦\ c\ •\ B\ t)\ a\ h\ =\ c\ •\ ouInputIncrement\ γ\ B\ a\ h}
\end{leancode}
\leansource{\texttt{ouInputIncrement\_\allowbreak{}map} \quad KineticLangevin/\allowbreak{}CanonicalSplitting.lean, lines 26–44}
\begin{leancode}
\leanline{0}{theorem\ ouInputIncrement\_map\ \{E\ F\ :\ Type*\}\ [NormedAddCommGroup\ E]\ [NormedSpace\ ℝ\ E]}
\leanline{4}{[CompleteSpace\ E]\ [NormedAddCommGroup\ F]\ [NormedSpace\ ℝ\ F]\ [CompleteSpace\ F]}
\leanline{4}{(L\ :\ E\ →L[ℝ]\ F)\ (γ\ :\ ℝ)\ (B\ :\ ℝ\ →\ E)\ (hB\ :\ Continuous\ B)\ (a\ h\ :\ ℝ)\ :}
\leanline{4}{ouInputIncrement\ γ\ (fun\ t\ ↦\ L\ (B\ t))\ a\ h\ =\ L\ (ouInputIncrement\ γ\ B\ a\ h)}
\end{leancode}
\leanlabel{How the Lean proof follows the argument}
\texttt{ou\allowbreak{}Input\allowbreak{}Increment\_\allowbreak{}smul} gives the formal counterpart of this step. The main intermediate facts are \texttt{hi} in that order. Named proof applications include \texttt{interval\allowbreak{}Integral.\allowbreak{}integral\_\allowbreak{}congr}; the adjacent proof shows their hypotheses and order of use.




\texttt{ou\allowbreak{}Input\allowbreak{}Increment\_\allowbreak{}map} gives the formal counterpart of this step. The main intermediate facts are \texttt{hc}, \texttt{hi} in that order. Named proof applications include \texttt{interval\allowbreak{}Integral.\allowbreak{}integral\_\allowbreak{}congr}; the adjacent proof shows their hypotheses and order of use.
\end{leanbox}

\section{\texorpdfstring{The continuous \(L^{2}\) Brownian input}{The continuous L² Brownian input}}


The input is the already constructed vector Brownian increment from time zero, evaluated at \(t_+=\max\{t,0\}\). This supplies a globally defined real-time function, anchored at zero. The proved continuity of vector Brownian increments in \(L^{2}\) and continuity of \(t\mapsto t_+\) prove continuity of this input. At zero the underlying increment is the difference of a Brownian value with itself, so it is exactly zero. The proof unfolds the source definitions to check this anchor.



Assemble the specified scalar Brownian coordinates into a vector and lift them onto the product initialization space. Bounded linearity transfers second-mean continuity and the zero initial value. This produces the precise continuous Hilbert-space input to which splitting convergence applies.

\begin{leanbox}{Lean: \texttt{canonicalVectorInput}, \texttt{continuous\_\allowbreak{}canonicalVectorInput}, \texttt{canonicalVectorInput\_\allowbreak{}zero}}
\leansource{\texttt{canonicalVectorInput} \quad KineticLangevin/\allowbreak{}CanonicalSplitting.lean, lines 49–51}
\begin{leancode}
\leanline{0}{def\ canonicalVectorInput\ (t\ :\ ℝ)\ :}
\leanline{4}{Lp\ (EuclideanSpace\ ℝ\ ι)\ 2\ (ULA.CanonicalBrownian.vectorProbability\ ι)\ :=}
\leanline{2}{vectorBrownianIncrementLp\ ι\ 0\ t.toNNReal}
\end{leancode}
\leansource{\texttt{continuous\_\allowbreak{}canonicalVectorInput} \quad KineticLangevin/\allowbreak{}CanonicalSplitting.lean, lines 53–54}
\begin{leancode}
\leanline{0}{theorem\ continuous\_canonicalVectorInput\ :\ Continuous\ (canonicalVectorInput\ ι)}
\end{leancode}
\leansource{\texttt{canonicalVectorInput\_\allowbreak{}zero} \quad KineticLangevin/\allowbreak{}CanonicalSplitting.lean, lines 56–57}
\begin{leancode}
\leanline{0}{@[simp]\ theorem\ canonicalVectorInput\_zero\ :\ canonicalVectorInput\ ι\ 0=0}
\end{leancode}
\leanlabel{How the Lean proof follows the argument}
\texttt{canonical\allowbreak{}Vector\allowbreak{}Input} names the mathematical object defined above. Its defining expression appears in the adjacent statement; subsequent identities may unfold this definition directly.




\texttt{continuous\_\allowbreak{}canonical\allowbreak{}Vector\allowbreak{}Input} gives the formal counterpart of this step.




\texttt{canonical\allowbreak{}Vector\allowbreak{}Input\_\allowbreak{}zero} gives the formal counterpart of this step. Its remaining equalities are discharged by the simplifications displayed in the proof.
\end{leanbox}


Earlier results used here: \cite[Section~3.8.17]{LangevinLean2026ULASC}.

\section{\texorpdfstring{Identify input increments and OU segment responses}{Identify input increments and OU segment responses}}


The difference between the anchored values at \(a+t\) and \(a\) cancels their common time-zero value, leaving the specified Brownian increment on \((a,a+t]\). The independently proved segment convolution identity gives \[Z_{a,t}=\Delta B_{a,t}-\gamma\int_0^t e^{-\gamma(t-s)}\Delta B_{a,s}\,ds.\] This is exactly the filtered Ornstein–Uhlenbeck increment used in the splitting theorem. The argument rewrites the endpoint increment, compares the interval integrands, and uses nonnegativity of the integration time to remove its real-to-nonnegative-time clamp. Scaling by any \(c\), including \(\sqrt{2\gamma}\), follows from the linearity identity. The equality is of \(L^{2}\) variables, stronger than merely matching marginal Gaussian laws.



The increment of the full Brownian input over a grid interval equals the corresponding segment input. Its filtered response equals the splitting thermostat innovation. These identities are proved from the source kernels, so the discretization uses the same forcing as the limiting solution.

\begin{leanbox}{Lean: \texttt{canonicalVectorInput\_\allowbreak{}increment}, \texttt{canonicalVectorInput\_\allowbreak{}ouIncrement}, \texttt{canonicalVectorInput\_\allowbreak{}scaled\_\allowbreak{}ouIncrement}}
\leansource{\texttt{canonicalVectorInput\_\allowbreak{}increment} \quad KineticLangevin/\allowbreak{}CanonicalSplitting.lean, lines 61–66}
\begin{leancode}
\leanline{0}{theorem\ canonicalVectorInput\_increment\ (a\ t\ :\ ℝ≥0)\ :}
\leanline{4}{canonicalVectorInput\ ι\ ((a+t\ :\ ℝ≥0)\ :\ ℝ)-canonicalVectorInput\ ι\ (a\ :\ ℝ)\ =}
\leanline{6}{vectorBrownianIncrementLp\ ι\ a\ t}
\end{leancode}
\leansource{\texttt{canonicalVectorInput\_\allowbreak{}ouIncrement} \quad KineticLangevin/\allowbreak{}CanonicalSplitting.lean, lines 70–85}
\begin{leancode}
\leanline{0}{theorem\ canonicalVectorInput\_ouIncrement\ (γ\ :\ ℝ)\ (hγ\ :\ 0≤γ)\ (a\ t\ :\ ℝ≥0)\ :}
\leanline{4}{ouInputIncrement\ γ\ (canonicalVectorInput\ ι)\ a\ t\ =}
\leanline{6}{vectorSegmentLp\ ι\ (psi\ γ\ 0)\ (continuous\_psi\ γ\ 0)\ a\ t}
\end{leancode}
\leansource{\texttt{canonicalVectorInput\_\allowbreak{}scaled\_\allowbreak{}ouIncrement} \quad KineticLangevin/\allowbreak{}CanonicalSplitting.lean, lines 89–92}
\begin{leancode}
\leanline{0}{theorem\ canonicalVectorInput\_scaled\_ouIncrement\ (γ\ c\ :\ ℝ)\ (hγ\ :\ 0≤γ)\ (a\ t\ :\ ℝ≥0)\ :}
\leanline{4}{ouInputIncrement\ γ\ (fun\ s\ ↦\ c\ •\ canonicalVectorInput\ ι\ s)\ a\ t\ =}
\leanline{6}{c\ •\ vectorSegmentLp\ ι\ (psi\ γ\ 0)\ (continuous\_psi\ γ\ 0)\ a\ t}
\end{leancode}
\leanlabel{How the Lean proof follows the argument}
\texttt{canonical\allowbreak{}Vector\allowbreak{}Input\_\allowbreak{}increment} gives the formal counterpart of this step. Its remaining equalities are discharged by the simplifications displayed in the proof.




\texttt{canonical\allowbreak{}Vector\allowbreak{}Input\_\allowbreak{}ou\allowbreak{}Increment} gives the formal counterpart of this step. The main intermediate facts are \texttt{he} in that order. Named proof applications include \texttt{interval\allowbreak{}Integral.\allowbreak{}integral\_\allowbreak{}congr}; the adjacent proof shows their hypotheses and order of use.




\texttt{canonical\allowbreak{}Vector\allowbreak{}Input\_\allowbreak{}scaled\_\allowbreak{}ou\allowbreak{}Increment} gives the formal counterpart of this step.
\end{leanbox}

\section{\texorpdfstring{The entire constructed innovation chain}{The entire constructed innovation chain}}


On a mesh of nonnegative size \(h\), the continuous-input chain uses the increment \(N_{kh,h}\). The preceding identity replaces this by \(cZ_{kh,h}\) for every grid index. The explicit-innovation recursion therefore agrees exactly with the continuous-input recursion. First, the argument uses the proved equality between the two chain representations and then substitutes the OU increment identity at each nonnegative grid time.



The entire stream has the product of its explicit thermal Gaussian laws. Equality is established for the source variables, including their dependence across response coordinates within an interval. This permits the deterministic splitting sequence to be identified with the intended random chain.

\begin{leanbox}{Lean: \texttt{canonicalVectorInput\_\allowbreak{}chain}}
\leansource{\texttt{canonicalVectorInput\_\allowbreak{}chain} \quad KineticLangevin/\allowbreak{}CanonicalSplitting.lean, lines 96–109}
\begin{leancode}
\leanline{0}{theorem\ canonicalVectorInput\_chain\ (g\ :}
\leanline{6}{Lp\ (EuclideanSpace\ ℝ\ ι)\ 2\ (ULA.CanonicalBrownian.vectorProbability\ ι)\ →}
\leanline{6}{Lp\ (EuclideanSpace\ ℝ\ ι)\ 2\ (ULA.CanonicalBrownian.vectorProbability\ ι))}
\leanline{4}{(γ\ c\ :\ ℝ)\ (hγ\ :\ 0≤γ)\ (h\ :\ ℝ≥0)}
\leanline{4}{(initial\ :\ Lp\ (EuclideanSpace\ ℝ\ ι)\ 2\ (ULA.CanonicalBrownian.vectorProbability\ ι)\ ×}
\leanline{6}{Lp\ (EuclideanSpace\ ℝ\ ι)\ 2\ (ULA.CanonicalBrownian.vectorProbability\ ι))\ (n\ :\ ℕ)\ :}
\leanline{4}{kineticSplitChain\ g\ γ\ h\ initial\ (fun\ s\ ↦\ c\ •\ canonicalVectorInput\ ι\ s)\ n\ =}
\leanline{6}{kineticSplitNoiseChain\ g\ γ\ h\ initial}
\leanline{8}{(fun\ k\ ↦\ c\ •\ vectorSegmentLp\ ι\ (psi\ γ\ 0)\ (continuous\_psi\ γ\ 0)\ ((k\ :\ ℝ≥0)*h)\ h)\ n}
\end{leancode}
\leanlabel{How the Lean proof follows the argument}
\texttt{canonical\allowbreak{}Vector\allowbreak{}Input\_\allowbreak{}chain} gives the formal counterpart of this step. Its remaining equalities are discharged by the simplifications displayed in the proof.
\end{leanbox}


Earlier results used here: \cite[Section~3.8.17]{LangevinLean2026ULASC}.

\section{\texorpdfstring{Conclude \(L^{2}\) convergence on the specified probability source}{Conclude L² convergence on the specified probability source}}


Apply the generic Banach splitting theorem to the continuous constructed input \(cB\), its proved zero initial value, and the continuous kinetic integral solution \(Y\) driven by that same input. Rewrite the resulting continuous-input chain by the proved constructed innovation-chain identity. This yields convergence of the explicit chain with Gaussian convolution innovations on the mesh \(T/(n+1)\) to \(Y_T\). The argument checks the nonnegative-time casts in the mesh and grid endpoints. The integral-solution premise is the equation constructed elsewhere by Picard existence; no stationarity or numerical convergence is assumed as a hypothesis.



Apply the pathwise splitting convergence theorem in the Hilbert space of square-integrable variables. All continuity and increment premises have already been verified on the specified probability source. The conclusion is second-mean convergence of the sampler splitting to the constructed kinetic solution.

\begin{leanbox}{Lean: \texttt{canonicalVectorInput\_\allowbreak{}chain\_\allowbreak{}tendsto}}
\leansource{\texttt{canonicalVectorInput\_\allowbreak{}chain\_\allowbreak{}tendsto} \quad KineticLangevin/\allowbreak{}CanonicalSplitting.lean, lines 113–134}
\begin{leancode}
\leanline{0}{theorem\ canonicalVectorInput\_chain\_tendsto\ (g\ :}
\leanline{6}{Lp\ (EuclideanSpace\ ℝ\ ι)\ 2\ (ULA.CanonicalBrownian.vectorProbability\ ι)\ →}
\leanline{6}{Lp\ (EuclideanSpace\ ℝ\ ι)\ 2\ (ULA.CanonicalBrownian.vectorProbability\ ι))}
\leanline{4}{(M\ γ\ :\ ℝ≥0)\ (hg\ :\ LipschitzWith\ M\ g)\ (c\ :\ ℝ)}
\leanline{4}{(initial\ :\ Lp\ (EuclideanSpace\ ℝ\ ι)\ 2\ (ULA.CanonicalBrownian.vectorProbability\ ι)\ ×}
\leanline{6}{Lp\ (EuclideanSpace\ ℝ\ ι)\ 2\ (ULA.CanonicalBrownian.vectorProbability\ ι))}
\leanline{4}{(Y\ :\ ℝ\ →\ Lp\ (EuclideanSpace\ ℝ\ ι)\ 2\ (ULA.CanonicalBrownian.vectorProbability\ ι)\ ×}
\leanline{6}{Lp\ (EuclideanSpace\ ℝ\ ι)\ 2\ (ULA.CanonicalBrownian.vectorProbability\ ι))}
\leanline{4}{(hY\ :\ Continuous\ Y)\ (T\ :\ ℝ≥0)}
\leanline{4}{(heY\ :\ ∀\ t\ ∈\ Icc\ 0\ (T\ :\ ℝ),\ Y\ t-initial\ =}
\leanline{6}{-(∫\ s\ in\ 0..t,kineticNegativeDrift\ g\ γ\ 1\ (Y\ s))+(c\ •\ canonicalVectorInput\ ι\ t,0))\ :}
\leanline{4}{Tendsto\ (fun\ n\ :\ ℕ\ ↦\ kineticSplitNoiseChain\ g\ γ}
\leanline{6}{((T/((n+1\ :\ ℕ)\ :\ ℝ≥0)\ :\ ℝ≥0)\ :\ ℝ)\ initial}
\leanline{6}{(fun\ k\ ↦\ c\ •\ vectorSegmentLp\ ι\ (psi\ γ\ 0)\ (continuous\_psi\ γ\ 0)}
\leanline{8}{((k\ :\ ℝ≥0)*(T/((n+1\ :\ ℕ)\ :\ ℝ≥0)))\ (T/((n+1\ :\ ℕ)\ :\ ℝ≥0)))\ (n+1))}
\leanline{6}{atTop\ ({\leanfallback 𝓝}\ (Y\ T))}
\end{leancode}
\leanlabel{How the Lean proof follows the argument}
\texttt{canonical\allowbreak{}Vector\allowbreak{}Input\_\allowbreak{}chain\_\allowbreak{}tendsto} gives the formal counterpart of this step. The main intermediate facts are \texttt{he} in that order. Its remaining equalities are discharged by the simplifications displayed in the proof.
\end{leanbox}


Earlier results used here: \cite[Section~3.8.17]{LangevinLean2026ULASC}.

\begin{thebibliography}{ULA-LSI}
\bibitem[ULA-SC]{LangevinLean2026ULASC}
LangevinLean Project. \emph{Unadjusted Langevin under Strong Convexity}. LL-01, 2026. Draft report, content revision d50c0f64815b. \url{https://odakyildiz.github.io/LangevinLean/docs/pdf/UnadjustedLangevin/FullProof.pdf}.
\end{thebibliography}

\end{document}
